%% kapa: exam=1 round=25 year=2014 period=2 form=A subject=appraisal_law subject_ko=부동산관계법규 seq=37 number=37
%% source: Q-Net 감정평가사 기출문제 (cst00309) artlSeq=5101207 file=extracted/1008/5101207/2014년 제25회 감정평가사 1차시험 기출문제/2교시(A형).hwp
%% conversion: kordoc 4.24.51 -> tools/gen1.py (auto); review=proofread-sonnet; flags=none
\begin{question}{37}
부동산등기법령상 부동산의 표시에 관한 변경등기에 대한 설명으로 옳은 것은?
\begin{choices}
  \item 토지의 지번에 변경이 있는 경우 토지 소유권의 등기명의인은 그 사실이 있는 때부터 1개월 이내에 변경등기를 신청하여야 한다.
  \item 등기관이 지적소관청으로부터 토지의 표시와 지적공부가 일치하지 아니하다는 사실을 통지받은 경우에는 통지받은 날로부터 1개월 이내에 직권으로 표시변경의 등기를 하여야 한다.
  \item 등기관이 직권으로 토지의 표시변경등기를 하였을 때에 등기명의인이 2인 이상인 경우에는 변경등기의 사실을 그 모두에게 통지하여야 한다.
  \item 전세권에 관한 등기가 있는 건물에 관하여는 합병의 등기를 할 수 없다.
  \item 구분건물로서 그 대지권의 변경이 있는 경우 구분건물의 소유권의 등기명의인은 1동의 건물에 속하는 다른 구분건물의 소유권의 등기명의인을 대위하여 변경등기를 신청할 수 없다.
\end{choices}
\end{question}
\answer{1}
